\chapter{Data Contracts and Docker Deployment}
\label{appendix:schema}

\section{Core Pydantic Schema Definitions}
The primary state contracts are implemented as immutable Pydantic models:

\begin{lstlisting}[language=Python]
class IncidentVector(BaseModel):
    event_type: IncidentType
    affected_node_ids: tuple[str, ...]
    severity: SeverityLevel
    duration_minutes: int
    description: str
    lane_closure_ratio: float | None = None
    demand_multiplier: float | None = None
    signal_plan_delta: dict[str, Any] | None = None

class CandidateAction(BaseModel):
    node_id: str
    green_time_ratio: float
    action_kind: str = "candidate_action"
    executable: bool = False
    requires_operator_approval: bool = True
    applied_by_system: bool = False

class WhatIfJobResult(BaseModel):
    job_id: str
    trace_id: str
    baseline_summary: dict[str, Any]
    scenario_summary: dict[str, Any]
    safety_iterations: int
    recommended_action: CandidateAction | None = None
    candidate_action: CandidateAction | None = None
    citations: list[dict[str, Any]] = []
    network_impact: dict[str, Any] | None = None
    audit_record: dict[str, Any]
\end{lstlisting}

\section{Docker Infrastructure}
The microservices stack is provisioned via Docker Compose:
\begin{itemize}[leftmargin=1.5cm, label=$\bullet$]
  \item \textbf{\texttt{gateway}:} FastAPI web server running Uvicorn on port 8000.
  \item \textbf{\texttt{worker}:} Distributed Celery task consumers processing LangGraph state machines.
  \item \textbf{\texttt{redis}:} Task broker and pub/sub message bus on port 6379.
  \item \textbf{\texttt{timescaledb}:} Time-series traffic database on port 5432.
  \item \textbf{\texttt{qdrant}:} Vector database on port 6333.
\end{itemize}
